\chapter{ÁREA DE APLICACIÓN}
\label{cap:area-aplicacion}

El presente capítulo describe el enfoque metodológico adoptado para el
desarrollo del sistema inteligente de entrenamiento en gimnasios con realidad
aumentada e inteligencia artificial. El propósito es garantizar un proceso de
construcción ordenado, iterativo y orientado a la calidad, acorde a las
necesidades específicas del desarrollo móvil. Para ello, se selecciona la
metodología ágil Mobile-D, diseñada especialmente para proyectos móviles con
ciclos de producción rápidos, entregas frecuentes y constante
retroalimentación.

% =====================================================================
\section{Antecedentes}
\label{sec:aa-antecedentes}

\subsection{Enfoque metodológico}

La metodología es el conjunto de procedimientos, técnicas y estrategias que
orienta la ejecución de un proyecto de investigación o desarrollo tecnológico.
En el presente trabajo, el enfoque metodológico define la manera en que se
estructura el proceso de diseñar, desarrollar y evaluar el sistema. El proyecto
adopta un enfoque cuantitativo-aplicado porque se medirá métricas objetivas del
sistema inteligente, como también se realizará la creación de una solución
tecnológica concreta mediante el uso de técnicas de ingeniería de software,
inteligencia artificial y realidad aumentada. En este caso, se emplean
fundamentos de biomecánica, visión por computadora, aprendizaje automático y
diseño de software para desarrollar una herramienta que facilite el uso
correcto de máquinas, la corrección técnica de ejercicios y la personalización
del entrenamiento.

% =====================================================================
\section{Metodología de desarrollo}
\label{sec:aa-metodologia}

La metodología Mobile-D fue seleccionada para guiar el proceso de desarrollo
del sistema debido a que está específicamente diseñada para proyectos de
software orientados a dispositivos móviles, caracterizados por ciclos de
producción rápidos, restricciones de hardware y cambios frecuentes en los
requerimientos.

Mobile-D es una metodología ágil diseñada para el desarrollo de aplicaciones
móviles, que requiere ciclos de entrega cortos, validación continua, pruebas en
dispositivos reales y flexibilidad ante cambios de interfaz o funcionalidad. El
presente proyecto integra módulos de guía interactiva mediante imágenes,
animaciones o realidad aumentada, generación de rutinas personalizadas mediante
inteligencia artificial, herramientas de seguimiento físico y una interfaz
altamente intuitiva.

Mobile-D permite abordar de forma progresiva y controlada, los elementos
principales del proyecto:

\begin{itemize}
    \item La guía visual sobre el uso correcto de máquinas y la ejecución de
    ejercicios, que exige pruebas de usabilidad desde las primeras etapas.
    \item La generación de rutinas personalizadas mediante algoritmos de IA,
    cuya precisión y utilidad deben ajustarse iterativamente.
    \item Los módulos de planificación y seguimiento, que requieren
    refinamientos sucesivos para asegurar claridad, funcionalidad y coherencia
    en la experiencia del usuario.
    \item La interfaz móvil, que debe ser evaluada continuamente para
    garantizar facilidad de uso, accesibilidad y entendimiento para usuarios
    principiantes.
\end{itemize}

% =====================================================================
\section{Aplicación de la Metodología Mobile-D}
\label{sec:aa-aplicacion-mobiled}

La metodología Mobile-D se implementa en el presente proyecto mediante cinco
fases que permiten estructurar el desarrollo de la aplicación de manera
incremental, verificable y orientada a la calidad. A continuación, se detalla
la aplicación específica de cada fase al sistema inteligente de entrenamiento.

% ---------------------------------------------------------------------
\subsection{Fase de Exploración}

En esta etapa se definen las bases conceptuales, técnicas y organizativas del
presente proyecto.

\titulillo{Módulos principales del sistema}

El sistema propuesto está compuesto por los siguientes módulos esenciales:

\begin{itemize}
    \item Registro y autenticación de usuarios.
    \item Guía detallada del uso correcto de máquinas de gimnasio.
    \item Ejecución técnica correcta de ejercicios específicos.
    \item Rutinas personalizadas mediante inteligencia artificial.
    \item Registro nutricional básico (calorías diarias).
    \item Seguimiento del progreso físico (peso, IMC, medidas).
    \item Reportes estadísticos y visualización del avance.
\end{itemize}

\titulillo{Requerimientos de sistema}

\textit{Requerimientos Funcionales}

\begin{itemize}
    \item RF1: Registrar y autenticar usuarios mediante correo y contraseña.
    \item RF2: Mostrar información técnica y visual de cada máquina de gimnasio.
    \item RF3: Generar rutinas personalizadas mediante IA según edad, peso, nivel
    y objetivo.
    \item RF4: Registrar peso, estatura, IMC y progreso físico.
    \item RF5: Registrar ingesta calórica diaria del usuario.
    \item RF6: Visualizar reportes y estadísticas del progreso físico.
\end{itemize}

\textit{Requerimientos No Funcionales}

\begin{itemize}
    \item RNF1: La aplicación debe ejecutarse en Android versión 8.0 o superior.
    \item RNF2: La interfaz debe ser intuitiva, rápida y adaptable a diferentes
    pantallas.
    \item RNF3: La sincronización de datos debe realizarse en tiempo real
    mediante Firebase.
\end{itemize}

\titulillo{Selección de Tecnologías}

\textit{Frontend}

\begin{itemize}
    \item \textbf{React Native:} Es un framework multiplataforma, que permite
    desarrollar aplicaciones móviles para dispositivos Android.
    \item \textbf{JavaScript:} Lenguaje de programación principal para el
    desarrollo de las interfaces y la lógica del lado del cliente.
    \item \textbf{Figma:} Herramienta de diseño utilizada para crear prototipos
    interactivos, pantallas iniciales, facilitando la validación visual del
    producto.
\end{itemize}

\textit{Backend}

\begin{itemize}
    \item \textbf{Firebase Authentication:} Servicio que permite gestionar el
    registro, inicio de sesión y recuperación de credenciales de forma segura.
    \item \textbf{Firebase Firestore:} Base de datos NoSQL en la nube que
    proporciona sincronización en tiempo real y escalabilidad para manejar
    información de usuarios, ejercicios, rutinas y progreso.
    \item \textbf{Firebase Storage:} Servicio destinado al almacenamiento de
    imágenes, animaciones y videos que acompañan la guía técnica del sistema.
    \item \textbf{TensorFlow Lite:} Biblioteca optimizada para ejecutar modelos
    de inteligencia artificial directamente en el dispositivo móvil, utilizada
    para generar rutinas personalizadas basadas en parámetros del usuario.
\end{itemize}

% ---------------------------------------------------------------------
\subsection{Fase de Inicialización}

En esta fase se prepara el entorno de desarrollo, la planificación inicial y la
definición de la arquitectura del sistema. También se establecen los
lineamientos de codificación y los mecanismos de control de versiones (Git).

\titulillo{Configuración del entorno}

Las principales actividades a realizar para la preparación del ambiente, son:

\begin{itemize}
    \item Instalación de Node.js.
    \item Instalación de Expo CLI.
    \item Creación de proyecto en Firebase y configuración de los servicios de
    autenticación, Firestore y Storage.
    \item Integración de la librería de TensorFlow Lite.
    \item Creación del repositorio en GitHub y definición de las ramas
    principales.
\end{itemize}

\titulillo{Arquitectura del sistema}

La arquitectura definida para el desarrollo de la aplicación es el de
cliente-servidor, donde la aplicación móvil actúa como cliente y los servicios
en la nube proporcionan autenticación, almacenamiento de datos, persistencia
multimedia y ejecución de algoritmos de IA.

La arquitectura cliente servidor permite garantizar escalabilidad,
sincronización en tiempo real y disponibilidad permanente del servicio.

\titulillo{Estructura de la arquitectura}

La estructura de la arquitectura del sistema se encuentra basada en capas:

\paragraph{Capa de Presentación}

Tecnologías: React Native + Expo, JavaScript, componentes UI.

\textit{Responsabilidades:}

\begin{itemize}
    \item Generar la interfaz visual de la aplicación.
    \item Gestionar la interacción del usuario.
    \item Capturar entradas como datos personales, selección de máquinas,
    rutinas y registro nutricional.
    \item Comunicar solicitudes al backend y mostrar respuestas.
\end{itemize}

\textit{Componentes clave:}

\begin{itemize}
    \item Módulo de registro y autenticación.
    \item Módulo de guía de máquinas y ejercicios.
    \item Módulo de rutinas personalizadas.
    \item Módulo de nutrición.
    \item Módulo de progreso y reportes.
\end{itemize}

\paragraph{Capa de Lógica de Negocio}

Esta se encuentra dentro de la app y está conectada a servicios externos.

\textit{Responsabilidades:}

\begin{itemize}
    \item Manejar reglas del sistema (cálculo de IMC, validaciones).
    \item Procesar datos antes de enviarlos a Firebase.
    \item Integrar y ejecutar modelos TensorFlow Lite para personalizar
    rutinas.
    \item Gestionar filtros, búsqueda de máquinas y procesamiento de
    ejercicios.
\end{itemize}

\textit{Componentes de negocio:}

\begin{itemize}
    \item Controlador de rutinas (IA).
    \item Controlador de ejercicios y máquinas.
    \item Controlador de progreso físico.
    \item Controlador nutricional.
\end{itemize}

\paragraph{Capa de Servicios}

Esta capa proporciona servicios esenciales para la operación del sistema:

\textit{Firebase Authentication:}

\begin{itemize}
    \item Registro de usuarios.
    \item Inicio de sesión.
    \item Recuperación de contraseñas.
    \item Seguridad y manejo de tokens.
\end{itemize}

\textit{Firebase Firestore:}

\begin{itemize}
    \item Base de datos NoSQL en tiempo real.
    \item Almacenamiento de rutinas, registros, progreso y máquinas.
\end{itemize}

\textit{Firebase Storage:}

\begin{itemize}
    \item Almacenamiento de imágenes, videos o animaciones educativas.
    \item Contenido para la guía técnica de máquinas.
\end{itemize}

\textit{Firebase Cloud Functions (opcional):}

\begin{itemize}
    \item Validación de datos.
    \item Preprocesamiento de peticiones.
    \item Actualizaciones automáticas de rutinas.
\end{itemize}

Estas funciones se integran a través de peticiones HTTP o SDK en tiempo real.

\paragraph{Capa de Inteligencia Artificial}

Tecnología: TensorFlow Lite.

\textit{Responsabilidades:}

\begin{itemize}
    \item Ejecutar modelos de recomendación ligera para rutinas.
    \item Ajustar parámetros según peso, nivel, objetivo y progreso del
    usuario.
    \item Proveer resultados inmediatos sin depender del servidor.
\end{itemize}

Este esquema reduce latencia, permite personalización offline y mejora la
experiencia del usuario.

% ---------------------------------------------------------------------
\subsection{Fase de Producción}

Esta es la etapa central del desarrollo. Dentro de esta etapa se implementan
los módulos de IA, seguimiento físico y rutinas personalizadas. Esta fase de
producción se organiza en iteraciones cortas, cada una de las cuales agrega un
conjunto de funcionalidades al sistema.

\titulillo{\textit{Día de prueba}}

En este día se verifica que toda la plataforma técnica funcione correctamente
antes de empezar a desarrollar funcionalidades reales.

Se debe realizar las siguientes acciones:

\begin{itemize}
    \item \textbf{Verificación de la integración con Firebase}

    Se realizarán pruebas básicas con los servicios principales de Firebase
    para garantizar su correcto funcionamiento:
    \begin{itemize}
        \item \textbf{Authentication:} creación de usuarios de prueba,
        validación de inicio de sesión y manejo de errores.
        \item \textbf{Firestore:} creación de colecciones temporales, escritura
        y lectura de documentos, y verificación de permisos básicos.
        \item \textbf{Storage:} carga y visualización de archivos multimedia
        para garantizar el acceso correcto al módulo de almacenamiento.
    \end{itemize}
    Estas pruebas aseguran que la aplicación pueda gestionar información real
    durante la fase de producción.
\end{itemize}

\textbf{Validación del entorno móvil}

Se verifica la correcta compilación del proyecto tanto en un dispositivo físico
como en un emulador. Además, se comprueba la navegación entre pantallas
iniciales y la ejecución estable del entorno Expo.

\textbf{Prueba preliminar del modelo TensorFlow Lite}

Se verifica la carga del modelo de inteligencia artificial y su ejecución
básica en el dispositivo, con el fin de confirmar la compatibilidad entre
TensorFlow Lite y la plataforma móvil utilizada.

\textbf{Validación del control de versiones}

Se realiza una prueba de uso del repositorio GitHub, asegurando que se pueden
subir y descargar cambios correctamente, lo cual es indispensable para mantener
un desarrollo organizado.

\titulillo{\textit{Día de salida}}

En este día se cierra la fase de inicialización, aquí no se realizan pruebas
sino ajustes finales y organización para dejar el proyecto listo para comenzar
la fase de producción.

Para esto, se debe realizar:

\textbf{Organización y limpieza del proyecto}

Se reorganiza la estructura interna del proyecto siguiendo la arquitectura
definida, agrupando componentes, servicios, pantallas y modelos de datos en
carpetas específicas. Asimismo, se eliminaron archivos temporales o de prueba
utilizados en fases previas.

\textbf{Configuración de variables y llaves de entorno}

Se definieron archivos seguros para almacenar las credenciales de Firebase y
otras configuraciones sensibles, evitando que estas sean expuestas en el
repositorio de código.

\textbf{Creación de ramas base de desarrollo en GitHub}

Se estableció una estructura de ramas para gestionar el desarrollo:

\begin{itemize}
    \item \texttt{main}: versión estable.
    \item \texttt{development}: versión en progreso.
    \item \texttt{feature/nombre-del-módulo}: para cada nueva funcionalidad.
\end{itemize}

Esta estructura permite un control ordenado de versiones.

% ---------------------------------------------------------------------
\subsection{Fase de Estabilización}

Dentro de esta fase se integran y verifican los componentes desarrollados,
asegurando consistencia, rendimiento y compatibilidad entre módulos. Las
pruebas de usabilidad y validación funcional se intensifican.

Las actividades a realizarse en esta fase son:

\begin{itemize}
    \item Refactorizar código para mejorar legibilidad.
    \item Optimizar carga de imágenes y videos.
    \item Ajustar tiempos de respuesta del sistema.
    \item Revisar el diseño visual para mejorar usabilidad.
    \item Corregir detalles detectados durante las pruebas en dispositivos
    reales.
\end{itemize}

% ---------------------------------------------------------------------
\subsection{Fase de Pruebas del sistema}

Se ejecutan pruebas de aceptación, evaluación del modelo de IA, pruebas en
dispositivos reales y corrección final de errores. El objetivo es garantizar
que el sistema sea estable y esté listo para su uso práctico.

Esto mediante las siguientes pruebas:

\begin{itemize}
    \item \textbf{Pruebas Funcionales:} Las pruebas funcionales son un tipo de
    verificación que evalúa si el sistema cumple correctamente con los
    requerimientos establecidos, las pruebas funcionales se basan en técnicas
    de caja negra, donde el evaluador verifica funciones, escenarios y
    respuestas del sistema desde la perspectiva del usuario final (IEEE, 1990).
    \item \textbf{Pruebas de Usabilidad:} Las pruebas de usabilidad se enfocan
    en analizar qué tan fácil, eficiente y satisfactoria es la interacción de un
    usuario con la interfaz del sistema. Se fundamenta en principios como
    facilidad de aprendizaje, eficiencia, reducción de errores y satisfacción
    del usuario (Nielsen, 1993).
    \item \textbf{Pruebas de Aceptación:} Las pruebas de aceptación validan que
    el sistema completo cumpla con los requerimientos acordados, estas pruebas
    representan la etapa final de validación del producto e incluyen la
    ejecución de escenarios reales que reflejan el uso cotidiano de la
    aplicación (Pressman, 2010).
\end{itemize}

% =====================================================================
\section{Roles}
\label{sec:aa-roles}

Dentro de la metodología Mobile-D, los roles cumplen la función de organizar
las responsabilidades durante el desarrollo del proyecto, aun cuando éste sea
desarrollado por una sola persona.

A continuación, se describen los roles adoptados en el proyecto:

\begin{table}[H]
\centering
\caption{Roles establecidos por Mobile-D}
\label{tab:roles}
\begin{tabularx}{\textwidth}{L{3.6cm} Y L{3cm}}
\toprule
\textbf{Rol} & \textbf{Responsabilidades} & \textbf{Responsable} \\
\midrule
Product Owner & Define alcance, objetivos y prioridades del sistema & Erika
Chino \\
Arquitecto Mobile & Diseña la arquitectura y selecciona las tecnologías & Erika
Chino \\
Desarrollador móvil & Implementa módulos, interfaces y lógica de negocio &
Erika Chino \\
Diseñador UI/UX & Diseña prototipos, navegación y experiencia de usuario &
Erika Chino \\
QA Tester & Realiza pruebas funcionales, de usabilidad y rendimiento & Erika
Chino \\
\bottomrule
\end{tabularx}
\fuente{Elaboración propia.}
\end{table}

% =====================================================================
\section{Técnicas e instrumentos de recolección de información}
\label{sec:aa-tecnicas}

Para fundamentar el diseño y funcionamiento del sistema se utilizarán las
siguientes técnicas:

\begin{itemize}
    \item \textbf{Revisión documental:} Investigación de aplicaciones similares,
    biomecánica, técnica correcta de ejercicios y libros especializados.
    \item \textbf{Observación directa:} Registro del uso real de máquinas en
    gimnasios locales.
    \item \textbf{Entrevistas informales:} Conversaciones con usuarios
    principiantes para conocer sus dificultades reales en el gimnasio.
    \item \textbf{Prototipado:} Elaboración de pantallas en Figma para validar
    la experiencia de usuario antes del desarrollo.
\end{itemize}
